\section{多项式环}\label{sec:galois-1-4}
\subsection{构造}\label{sec:galois-1-4-1}
设 $A$ 是交换环，$A^{(\mathbb N)}$ 是以 $\mathbb N$ 为基的自由 $A$-模. 为固定记号，将 $A^{(\mathbb N)}$ 的典范基记为 $\{e_n,n\in\mathbb N\}$. 我们将利用指标集 $\mathbb N$ 的加法，赋予 $A^{(\mathbb N)}$ 一个环结构. 为此，只须将典范基上由下式给出的乘法规则作线性延拓：
\begin{align*}
e_n\cdot e_m:=e_{n+m},\qquad \forall n,m\in\mathbb N.
\end{align*}
将元素 $(a_n)_{n\in\mathbb N}\in A^{\mathbb N}$ 写成 $(a_n)_{n\in\mathbb N}=\sum_{n\in\mathbb N}a_ne_n$，便得到显式公式\footnote{译注：此处原文写作 $A^{\mathbb N}$，但构造所在的模是有限支撑序列组成的 $A^{(\mathbb N)}$. 对一般的 $A^{\mathbb N}$ 元素，上述和不是自由模中的有限和.}
\begin{align*}
(a_n)_{n\in\mathbb N}\cdot(b_n)_{n\in\mathbb N}:=(c_n)_{n\in\mathbb N},\qquad c_n=\sum_{p+q=n}a_pb_q.
\end{align*}
容易验证，这一运算是 $A$-双线性的（特别地，它对加法满足分配律），并且满足结合律，以 $e_0=(1,0,\ldots)$ 为单位元. 因而它在 $A^{(\mathbb N)}$ 上定义了一个环结构，而环态射 $A\longrightarrow A^{(\mathbb N)}$，$a\longmapsto ae_0=(a,0,0,\ldots)$ 使之成为一个 $A$-代数.

将元素 $e_1=(0,1,0,0,\ldots)$ 记作 $X$. 由定义有 $X^n=e_n$，故族 $\{1,X,X^2,\ldots\}$ 就是 $A^{(\mathbb N)}$ 的典范基. 因此，每个元素都唯一地写成 $(a_n)_{n\in\mathbb N}=\sum_{n\in\mathbb N}a_nX^n$，其中当 $n\gg0$ 时 $a_n=0$，而乘法写成
\begin{align*}
\left(\sum_n a_nX^n\right)\cdot\left(\sum_n b_nX^n\right)=\sum_n\left(\sum_{p+q=n}a_pb_q\right)X^n.
\end{align*}
\begin{definition*}
上述 $A$-代数记作 $A[X]$，称为 $A$ 上的（关于不定元 $X$ 的）\newterm{galois-polynomial-algebra}{多项式代数}{polynomial algebra}. 它的元素称为系数在 $A$ 中的（关于不定元 $X$ 的）\newterm{galois-polynomial}{多项式}{polynomial}.
\end{definition*}

\subsection{泛性质}\label{sec:galois-1-4-2}
多项式 $A$-代数 $A[X]$ 连同其元素 $X$ 由下述泛性质刻画：
\begin{proposition*}
对任意 $A$-代数 $B$（无论交换与否）和任意元素 $b\in B$，存在唯一的 $A$-代数态射 $\varphi\colon A[X]\longrightarrow B$ 使得 $\varphi(X)=b$. 换言之，对任意 $A$-代数 $B$，映射 $\varphi\longmapsto\varphi(X)$ 是双射
\begin{align*}
\operatorname{Hom}_{A\text{-}\mathrm{alg}}(A[X],B)\xrightarrow{\sim}B.
\end{align*}
其逆双射将元素 $b\in B$ 对应到``在 $b$ 处的求值态射''，即将 $f=\sum_n a_nX^n$ 映到 $f(b):=\sum_n a_nb^n$ 的态射 $A[X]\longrightarrow B$.
\end{proposition*}
\begin{proof}
唯一性：我们要求 $\varphi(X)=b$，且 $\varphi$ 与环结构相容，故没有其他选择，必须有 $\varphi(\sum_n a_nX^n)=\sum_n a_nb^n$.
存在性：由于 $(1,X,X^2,\ldots)$ 是 $A$-模 $A[X]$ 的一组基，表达式 $\varphi(\sum_n a_nX^n):=\sum_n a_nb^n$ 定义了一个 $A$-模态射. 容易验证它与乘法相容；实际上，$A[X]$ 的乘法正是为此而构造的.
\end{proof}
可以用下面这句话非正式地概括上述命题：给定一个 $A$-代数态射 $A[X]\longrightarrow B$，等价于给定 $X$ 的像 $b\in B$.
\begin{corollary*}
设 $M$ 是 $A$-模，$\chi$ 是 $A$-模 $M$ 的一个自同态. 则 $M$ 上存在唯一的 $A[X]$-模结构，它延拓原来的 $A$-模结构，并且使 $X$ 通过 $\chi$ 作用. 因而，给定一个 $A[X]$-模，等价于给定一个由 $A$-模及其一个自同态组成的有序对.
\end{corollary*}
\begin{proof}
事实上，若 $A\xrightarrow{\rho}\operatorname{End}_{\mathbb Z}(M)$ 是定义 $M$ 的 $A$-模结构的环态射，则由前一命题，存在唯一的 $A$-代数态射 $A[X]\xrightarrow{\widetilde\rho}\operatorname{End}_{\mathbb Z}(M)$，它延拓 $\rho$ 并将 $X$ 映到 $\chi$. 该态射定义了所述的 $A[X]$-模结构.
\end{proof}
\begin{remark*}
$A[X]$ 的泛性质与商的泛性质，给出了形如 $A[X]/(f)$（其中 $f\in A[X]$）的 $A$-代数的一个泛性质. 事实上，由商的泛性质，给定态射 $A[X]/(f)\longrightarrow B$，等价于给定核包含 $f$ 的态射 $A[X]\xrightarrow{\varphi}B$. 而给定 $\varphi$ 又等价于给定 $b=\varphi(X)$，且``核包含 $f$''等价于 $\varphi(f)=0$. 由于 $\varphi(f)=f(b)$，我们看到，给定态射 $A[X]/(f)\longrightarrow B$，就等价于给定满足 $f(b)=0$ 的元素 $b\in B$.

例如，给定一个环态射 $\mathbb Z[X]/(X^n)\longrightarrow B$，等价于给定 $B$ 中一个幂零指数不大于 $n$ 的幂零元. 给定一个从 $\mathbb Z[X]/(X^n-1)$ 出发的环态射，等价于给定一个 $n$ 次单位根.
\end{remark*}


\subsection{性质的传递}\label{sec:galois-1-4-3}
我们将看到，环 $A$ 的某些性质会传递给环 $A[X]$. 先回顾几个定义.
\begin{definition*}
设 $f=\sum_n a_nX^n$ 是非零多项式.
\begin{enumerate}
\item[i)] 它的\newterm{galois-polynomial-degree}{次数}{degree} $\deg(f)$ 是使 $a_n\ne0$ 的最大指标 $n$.
\item[ii)] 它的\newterm{galois-leading-coefficient}{首项系数}{leading coefficient}是 $a_{\deg(f)}$.
\item[iii)] 若 $a_{\deg(f)}=1$，则称 $f$ 为\newterm{galois-monic-polynomial}{首一多项式}{monic polynomial}.
\end{enumerate}
\end{definition*}
\begin{proposition*}
设 $f=\sum_n a_nX^n$ 与 $g=\sum_n b_nX^n$ 是 $A[X]$ 中的非零多项式. 写 $fg=\sum_n c_nX^n$. 则 $\deg(fg)\leq\deg(f)+\deg(g)$，且等号成立当且仅当 $a_{\deg(f)}b_{\deg(g)}\ne0$. 在此情形下，有 $c_{\deg(fg)}=a_{\deg(f)}b_{\deg(g)}$.
\end{proposition*}
\begin{proof}
乘法的定义表明，$fg=a_{\deg(f)}b_{\deg(g)}X^{\deg(f)+\deg(g)}+\text{（次数更低的项）}$.
\end{proof}
\begin{remark*}
为使关于乘积次数的等式对因子不加限制也成立，我们约定零多项式的次数为 $-\infty$，且这一符号满足次序关系 $-\infty<0$ 和加法关系：对任意 $n\in\mathbb N$，$-\infty+n=-\infty$.\footnote{译注：原文此处说``等式''，但上述约定只使涉及零多项式的次数表达式有意义，并不能消除首项系数乘积为零所导致的严格不等式. 其后的习题正要求给出严格不等式的例子.}
\end{remark*}
\begin{exercise*}
给出满足 $\deg(fg)<\deg(f)+\deg(g)$ 的多项式的例子.
\end{exercise*}
\begin{corollary*}
若 $A$ 是整环，则 $A[X]$ 也是整环，且 $A[X]^\times=A^\times$.
\end{corollary*}
\begin{proof}
等式 $\deg(fg)=\deg(f)+\deg(g)$ 蕴含 $f,g\ne0\Rightarrow fg\ne0$，即 $A[X]$ 是整环. 最后，若 $fg=1$，由于 $\deg(1)=0$，必有 $\deg(f)=\deg(g)=0$，故 $f,g\in A$，最终得到 $f,g\in A^\times$.
\end{proof}
\begin{example*}
有时 $A[X]^\times$ 会严格大于 $A^\times$. 设 $A=\mathbb Z/p^2\mathbb Z$. 则多项式 $1+\overline pX$ 在 $A[X]$ 中可逆，逆元为 $1-\overline pX$. 事实上，有
\footnote{译注：原文中间一式为 $1-\overline p^{\,2}X$，已照录；乘积展开应为 $1-\overline p^{\,2}X^2$. 因 $\overline p^{\,2}=0$，所述可逆性结论不受影响.}
\begin{align*}
(1+\overline pX)(1-\overline pX)=1-\overline p^{\,2}X=1.
\end{align*}
\end{example*}
\begin{proposition*}[{Euclid 除法}]
设 $f\in A[X]$ 是首一多项式（因而非零）. 则对任意非零的 $g\in A[X]$，存在唯一的有序对 $(q,r)\in A[X]^2$，使得 $\deg(r)<\deg(f)$ 且 $g=qf+r$.
\end{proposition*}
\begin{proof}
存在性. 对 $\delta:=\deg(g)-\deg(f)$ 作归纳. 注意，若 $\delta<0$，可取 $q=0$、$r=g$. 另一方面，若 $\delta\geq0$，考虑 $g':=g-b_{\deg(g)}X^\delta f$，其中 $b_{\deg(g)}$ 是 $g$ 的首项系数. 显然 $\delta'=\deg(g')-\deg(f)<\delta$，由归纳假设，存在 $q',r'$ 使 $g'=q'f+r'$. 因而 $g=(q'+b_{\deg(g)}X^\delta)f+r'$，这正是所要的分解.

唯一性. 若 $qf+r=q'f+r'$，则 $(q-q')f=r'-r$. 假设 $r'\ne r$. 则 $\deg(r'-r)<\deg(f)$，而由于 $f$ 首一，有 $\deg(qf-qf')=\deg(q-q')+\deg(f)$，矛盾. 因而 $r=r'$，于是 $qf=qf'$，最后，由于首一多项式 $f$ 不是零因子，得到 $q=q'$.\footnote{译注：原文此处两次写成 $qf'$，已照录；与前式 $(q-q')f=r'-r$ 相应，应分别为 $\deg(qf-q'f)$ 和 $qf=q'f$.}
\end{proof}
\begin{corollary*}
设 $f\in A[X]$ 是次数为 $d$ 的首一多项式. 则商 $A[X]/(f)$ 是自由 $A$-模，以 $\{\overline1,\overline X,\ldots,\overline X^{d-1}\}$ 为基，其中 $\overline X$ 表示 $X$ 在 $A/(f)$ 中的像.\footnote{译注：原文此处写 $A/(f)$，已照录；由上下文应为 $A[X]/(f)$.}
\end{corollary*}
\begin{proof}
设 $g\in A[X]$. 写 $g=qf+r$，其中 $\deg(r)<\deg(f)$. 因而存在 $c_i\in A$，使 $r=\sum_{i=0}^{d-1}c_iX^i$. 则 $g$ 在 $A[X]/(f)$ 中的像 $\overline g$ 为 $\overline g=\overline r=\sum_{i=0}^{d-1}c_i\overline X^i$，这说明命题中的族确实是生成族. 另一方面，假设 $\sum_{i=1}^d c_i\overline X^i=0$. 则 $r:=\sum_{i=1}^d c_iX^i\in(f)$，由 Euclid 除法的唯一性，有 $r=0$，故对所有 $i$ 都有 $c_i=0$，从而命题中的族确实是 $A$-模 $A[X]/(f)$ 的一组基.\footnote{译注：原文在线性无关性的论证中，两处求和指标均为 $i=1,\ldots,d$，已照录. 对所声明的基，应检验指标 $i=0,\ldots,d-1$ 的线性关系；按原文的指标，不能由除法唯一性推出 $r=0$.}
\end{proof}
\begin{remark*}
注意，$A[X]/(f)$ 作为 $A[X]$-模当然不是自由的，除非 $f=1$.
\end{remark*}
\begin{theorem*}[{Hilbert 基定理}]
若 $A$ 是 Noether 环，则 $A[X]$ 也是 Noether 环.
\end{theorem*}
\begin{proof}
设 $I$ 是 $A[X]$ 的一个理想. 我们要证明它有限生成. 作为一般原则，由第~\ref{sec:galois-1-3-7}~分目，只需证明 $I/J$ 是 $A[X]/J$ 中的有限生成理想，其中 $J\subset I$ 是我们自行选择的有限生成子理想.

特别地，当 $I$ 包含一个首一多项式 $f$ 时，可取 $J=(f)$. 前一推论告诉我们，$A/J$ 是有限生成 $A$-模，因而是 Noether 模，所以 $I/J$ 在 $A$ 上有限生成，进而在 $A[X]$ 上也有限生成.\footnote{译注：原文此处写 $A/J$，已照录；由前一推论，所指应为 $A[X]/J$.}

不过，$I$ 可能不包含任何首一多项式. 在这种情形下，考虑所有 $f\in I$ 的首项系数组成的集合 $K\subset A$. 显然，这是 $A$ 的一个理想（请验证），因此它由元素 $a_1,\ldots,a_r$ 在 $A$ 上生成，因为 $A$ 是 Noether 环. 对每个 $i=1,\ldots,r$，选取首项系数为 $a_i$ 的多项式 $f_i\in I$，并将 $f_1,\ldots,f_r$ 生成的 $A[X]$ 的理想记为 $J\subset I$. 我们将证明，$I$ 在 $A/J$ 中的像 $I/J$ 是有限生成 $A$-模；由第一段，这足以得出结论.\footnote{译注：原文此处也写 $A/J$，已照录，应理解为 $A[X]/J$. 此外，按本分目只对非零多项式定义首项系数的约定，$K$ 的定义还须补入 $0$ 才能成为理想；原文没有明说这一点，也没有单独处理 $I=0$ 的情形.}

为此，令 $d=\max\{\deg(f_1),\ldots,\deg(f_r)\}$. 只需证明
\begin{quote}
\hypertarget{galois-hilbert-star}{(*)} 每个多项式 $f\in I$ 都模 $J$ 同余于一个次数小于 $d$ 的多项式.
\end{quote}
事实上，若承认 \hyperlink{galois-hilbert-star}{(*)}，则 $I/J$ 是 $I\cap(A+AX+\cdots+AX^{d-1})$ 的像. 后者是 $A+AX+\cdots+AX^{d-1}$ 的 $A$-子模，因而在 Noether 环 $A$ 上有限生成. 因此，$I/J$ 本身也有限生成.

下面对 $\deg(f)$ 作归纳来证明 \hyperlink{galois-hilbert-star}{(*)}. 当 $\deg(f)<d$ 时，结论不言而喻. 现设 $\deg(f)\geq d$，并对 $i=1,\ldots,r$ 记 $\delta_i:=\deg(f)-\deg(f_i)>0$.\footnote{译注：原文写 $\delta_i>0$，已照录；由所给假设只能推出 $\delta_i\geq0$，归纳论证也只需要非负性.} 则 $f$ 的首项系数具有形式 $a_{\deg(f)}=c_1a_1+\cdots+c_ra_r$，其中 $c_1,\ldots,c_r\in A$. 于是，若令 $f':=f-\sum_{i=1}^r c_iX^{\delta_i}f_i$，就有 $f'\equiv f\pmod J$ 且 $\deg(f')<\deg(f)$. 对 $f'$ 应用归纳假设即得结论.
\end{proof}


% 原文：印刷页34--39，1.4.4--1.4.7.
\subsection{推广}\label{sec:galois-1-4-4}
在构造``一个未定元''的多项式环时，我们首先在 $A^{(\mathbb N)}$ 上定义了乘法；这一乘法依赖于为 $A^{(\mathbb N)}$ 的标准基编号的集合 $\mathbb N$ 上的加法. 可以尝试将 $(\mathbb N,+,0)$ 替换为任意结合幺半群 $(\mathcal N,+,0)$，从而推广这一过程. 记 $(e_\nu)_{\nu\in\mathcal N}$ 为 $A^{(\mathcal N)}$ 的标准基（于是 $(a_\nu)_{\nu\in\mathcal N}=\sum_{\nu\in\mathcal N}a_\nu e_\nu$）. 我们需要把在这组基上由下式定义的乘法作线性延拓：
\begin{align*}
e_\nu\cdot e_{\nu'}=e_{\nu+\nu'}\qquad(\nu,\nu'\in\mathcal N).
\end{align*}
展开便得到公式
\begin{align*}
(a_\nu)_{\nu\in\mathcal N}\cdot(b_\nu)_{\nu\in\mathcal N}
&:=(c_\nu)_{\nu\in\mathcal N},&
c_\nu&=\sum_{\mu+\rho=\nu}a_\mu b_\rho.
\end{align*}
这里的和确实只有有限多个非零项，因为我们考虑的是 $A^{(\mathcal N)}$，而不是 $A^{\mathcal N}$.

容易验证，这一运算是 $A$-双线性的（因而特别地，对加法满足分配律），并且满足结合律，以 $e_0$ 为单位元. 此外，映射 $A\longrightarrow A^{(\mathcal N)}$，$a\longmapsto ae_0$ 是环同态，因而使 $A^{(\mathcal N)}$ 成为一个 $A$-代数.
\begin{definition*}
这样定义的 $A$-代数称为幺半群 $\mathcal N$ 的 \newterm{galois-monoid-algebra}{幺半群代数}{monoid algebra}，通常记作 $A[\mathcal N]$. 当 $\mathcal N$ 是群时，也称为 $A$ 上的\newterm{galois-group-algebra}{群代数}{group algebra}.
\end{definition*}
如同``通常的''多项式一样，也可以用泛性质刻画 $A^{(\mathcal N)}$.
\begin{proposition*}
设 $B$ 是一个 $A$-代数，配备了幺半群同态 $\nu\longmapsto b_\nu$，$(\mathcal N,+,0)\longrightarrow(B,\cdot,1)$. 则存在唯一的 $A$-代数同态 $\varphi\colon A[\mathcal N]\longrightarrow B$，使得对每个 $\nu\in\mathcal N$，都有 $\varphi(e_\nu)=b_\nu$. 换言之，对任意 $A$-代数 $B$，映射 $\varphi\longmapsto(\nu\longmapsto\varphi(e_\nu))$ 是双射
\begin{align*}
\operatorname{Hom}_{A\text{-}\mathrm{alg}}(A[\mathcal N],B)
\xrightarrow{\sim}\operatorname{Hom}_{\mathrm{mono}}(\mathcal N,(B,\cdot)).
\end{align*}
\end{proposition*}
回顾一下，``$\nu\longmapsto b_\nu$ 是幺半群同态''意味着对所有 $\nu,\nu'\in\mathcal N$，在 $B$ 中都有等式 $b_\nu b_{\nu'}=b_{\nu+\nu'}$.
\begin{proof}
\emergencystretch=3em
由于族 $(e_\nu)_{\nu\in\mathcal N}$ 是 $A[\mathcal N]$ 的一组基，自由 $A$-模的泛性质保证存在唯一的 $A$-模同态 $\varphi\colon A[\mathcal N]\longrightarrow B$，将 $e_\nu$ 映为 $b_\nu$. 特别地，命题中所要求的同态若存在，必定唯一. 对于存在性，只需验证这个同态 $\varphi$ 与乘法相容. 对基 $(e_\nu)_{\nu\in\mathcal N}$ 中的元素，这是显然的，因为
\begin{align*}
\varphi(e_\nu e_{\nu'})=\varphi(e_{\nu+\nu'})=b_{\nu+\nu'}=b_\nu b_{\nu'}.
\end{align*}
由线性性，这对所有元素都成立.
\end{proof}
\begin{remark*}
上述构造在非交换情形中同样完全适用. 若从交换环 $A$ 和非交换幺半群 $\mathcal N$ 出发，就得到非交换 $A$-代数 $A[\mathcal N]$. 最经典的例子是 $\mathcal N$ 为有限群 $G$ 的情形. 代数 $A[G]$ 就是``群 $G$ 的群代数''.
\end{remark*}
下面考察这一构造的几个有趣特例.

\subsection{多元多项式}\label{sec:galois-1-4-5}
这是 $\mathcal N=\mathbb N^n$ 的例子. 因而元素 $\nu\in\mathbb N^n$ 是一个 $n$ 元组 $(\nu_1,\ldots,\nu_n)$，加法逐项定义为 $\nu+\nu'=(\nu_1+\nu'_1,\ldots,\nu_n+n'_n)$.\footnote{译注：原文最后一项印作 $\nu_n+n'_n$；按逐项相加的定义，末项中的 $n'_n$ 应为 $\nu'_n$.} 在 $A[\mathbb N^n]$ 中令
\begin{align*}
X_i:=e_{(0,\ldots,1,\ldots,0)},\qquad\text{其中 $1$ 位于第 $i$ 个位置.}
\end{align*}
于是对每个 $\nu=(\nu_1,\ldots,\nu_n)$，都有等式
\begin{align*}
e_\nu=X_1^{\nu_1}X_2^{\nu_2}\cdots X_n^{\nu_n}.
\end{align*}
因此，族 $(X_1^{\nu_1}\cdots X_n^{\nu_n})_{(\nu_1,\ldots,\nu_n)\in\mathbb N^n}$ 就是 $A[\mathbb N^n]$ 的标准基. 换言之，$A[\mathbb N^n]$ 的每个元素 $f$ 都唯一地写成
\begin{align*}
f=\sum_{(\nu_1,\ldots,\nu_n)}a_{\nu_1,\ldots,\nu_n}X_1^{\nu_1}\cdots X_n^{\nu_n}.
\end{align*}
为简化记号，通常记 $X^\nu:=X_1^{\nu_1}\cdots X_n^{\nu_n}$，于是 $f=\sum_{\nu\in\mathbb N^n}a_\nu X^\nu$. 乘法由下式给出：
\begin{align*}
\left(\sum_{\nu\in\mathbb N^n}a_\nu X^\nu\right)
\left(\sum_{\nu\in\mathbb N^n}b_\nu X^\nu\right)
=\sum_\nu\left(\sum_{\mu+\rho=\nu}a_\mu b_\rho\right)X^\nu.
\end{align*}
\begin{definition*}
代数 $A[\mathbb N^n]$ 也记作 $A[X_1,\ldots,X_n]$，称为``$n$ 个未定元 $X_1,\ldots,X_n$ 的多项式代数''. 其中的元素称为关于 $X_i$ 的多项式，而元素 $X^\nu=X_1^{\nu_1}X_2^{\nu_2}\cdots X_n^{\nu_n}$ 称为\newterm{galois-monomial}{单项式}{monomial}.
\end{definition*}
因此，单项式构成 $A[X_1,\ldots,X_n]$ 的一组基，每个多项式都是单项式的 $A$-线性组合. 单项式 $X^\nu$ 的\newterm{galois-total-degree}{总次数}{total degree}定义为整数 $|\nu|=\nu_1+\cdots+\nu_n$. 多项式 $f=\sum_\nu a_\nu X^\nu$ 的总次数是所有满足 $a_\nu\ne0$ 的单项式 $X^\nu$ 的次数的最大值.

用 $X_i$ 来表述 $A[\mathbb N^n]$ 满足的泛性质更为方便：
\begin{proposition*}
对任意配备了元素 $b_1,\ldots,b_n$ 的 $A$-代数 $B$，存在唯一的 $A$-代数同态 $\varphi\colon A[X_1,\ldots,X_n]\longrightarrow B$，使得对每个 $i$ 都有 $\varphi(X_i)=b_i$. 换言之，对任意 $A$-代数 $B$，映射 $\varphi\longmapsto(\varphi(X_1),\ldots,\varphi(X_n))$ 是双射
\begin{align*}
\operatorname{Hom}_{A\text{-}\mathrm{alg}}(A[X_1,\ldots,X_n],B)\xrightarrow{\sim}B^n.
\end{align*}
其逆双射将元素 $(b_1,\ldots,b_n)\in B^n$ 对应到``在 $(b_1,\ldots,b_n)$ 处求值''的同态，即将 $f=\sum_\nu a_\nu X^\nu$ 映为 $f(b_1,\ldots,b_n):=\sum_\nu a_\nu b_1^{\nu_1}\cdots b_n^{\nu_n}$ 的同态 $A[X_1,\ldots,X_n]\longrightarrow B$.
\end{proposition*}
\begin{proof}
可以直接证明这一点. 但更自然的办法是由 $A[\mathbb N^n]$ 的泛性质推出它，只需注意，给定 $n$ 个元素 $b_1,\ldots,b_n$，等价于通过关系 $b_\nu=b_1^{\nu_1}b_2^{\nu_2}\cdots b_n^{\nu_n}$ 给定一个幺半群同态 $(\mathbb N^n,+,0)\longrightarrow(B,\cdot,1)$，$\nu\longmapsto b_\nu$.
\end{proof}
可以把这个命题非正式地改述为：给定同态 $A[X_1,\ldots,X_n]\longrightarrow B$，就相当于给定各个 $X_i$ 的像 $b_1,\ldots,b_n$.
\begin{remark*}
设 $f_1,\ldots,f_r$ 是 $A[X_1,\ldots,X_n]$ 中的多项式. 与前面 $n=r=1$ 的情形类似，给定 $A$-代数同态 $A[X_1,\ldots,X_n]/(f_1,\ldots,f_r)\longrightarrow B$，等价于给定满足 $f_1(b_1,\ldots,b_n)=\cdots=f_r(b_1,\ldots,b_n)=0$ 的 $n$ 元组 $(b_1,\ldots,b_n)\in B^n$.
\end{remark*}
\begin{corollary*}
存在典范的 $A$-代数同构 $A[X_1,\ldots,X_n]\xrightarrow{\sim}A[X_1,\ldots,X_{n-1}][X_n]$.
\end{corollary*}
\begin{proof}
只需证明 $A$-代数 $A[X_1,\ldots,X_{n-1}][X_n]$ 满足与 $A[X_1,\ldots,X_n]$ 相同的泛性质. 为此，设 $B$ 是一个 $A$-代数，配备了 $n$ 个元素 $b_1,\ldots,b_n$. 由 $A[X_1,\ldots,X_{n-1}]$ 的泛性质，存在一个从 $A[X_1,\ldots,X_{n-1}]$ 到 $B$ 的 $A$-代数同态，将 $X_i$ 映为 $b_i$，其中 $i=1,\ldots,n-1$. 这使 $B$ 成为一个 $A[X_1,\ldots,X_{n-1}]$-代数. 接着，单变量多项式的泛性质给出一个 $A[X_1,\ldots,X_{n-1}]$-代数同态
\begin{align*}
\varphi\colon A[X_1,\ldots,X_{n-1}][X_n]\longrightarrow B,
\end{align*}
将 $X_n$ 映为 $b_n$. 因此，$\varphi$ 也是一个 $A$-代数同态，对每个 $i$ 都将 $X_i$ 映为 $b_i$. 再证明这样的同态是唯一的. 若 $\varphi'$ 是另一个这样的同态，由 $A[X_1,\ldots,X_{n-1}]$ 的泛性质，有 $\varphi|_{A[X_1,\ldots,X_{n-1}]}=\varphi'|_{A[X_1,\ldots,X_{n-1}]}$；再由单变量多项式的泛性质，便得 $\varphi=\varphi'$.
\end{proof}
\begin{corollary*}
若 $A$ 是整环（相应地，Noether 环），则 $A[X_1,\ldots,X_n]$ 也是整环（相应地，Noether 环）.
\end{corollary*}
\begin{proof}
借助上一个推论，利用归纳法可归结到已经处理过的单变量情形.
\end{proof}
\paragraph{应用.}
设 $\mathcal F(\mathbb C^n)$ 是所有函数 $\mathbb C^n\longrightarrow\mathbb C$ 组成的 $\mathbb C$-代数. 其中有坐标函数 $z_1,\ldots,z_n$. 将 $X_i$ 映为 $z_i$ 的 $\mathbb C$-代数同态 $\mathbb C[X_1,\ldots,X_n]\longrightarrow\mathcal F(\mathbb C^n)$ 的像，定义为 $\mathbb C^n$ 上的多项式函数代数 $\mathcal O(\mathbb C^n)$. 换言之，函数是多项式函数，是指它具有形式 $(z_1,\ldots,z_n)\longmapsto f(z_1,\ldots,z_n)$，其中 $f\in\mathbb C[X_1,\ldots,X_n]$. 注意，这个多项式由该函数唯一确定（练习）. 换言之，同态 $\mathbb C[X_1,\ldots,X_n]\longrightarrow\mathcal F(\mathbb C^n)$ 是单射，可以将 $\mathbb C[X_1,\ldots,X_n]=\mathcal O(\mathbb C^n)$ 等同起来. 注意，泛性质给出双射
\begin{align*}
\mathbb C^n\xrightarrow{\sim}\operatorname{Hom}_{\mathbb C\text{-}\mathrm{alg}}(\mathcal O(\mathbb C^n),\mathbb C),\qquad z\longmapsto(f\longmapsto f(z)),
\end{align*}
使我们能够将``空间'' $\mathbb C^n$ 的每个点，看作该空间的函数代数上的一个求值同态.

更一般地，设 $V\subset\mathbb C^n$ 是 $\mathbb C^n$ 的一个代数子集，$\mathcal F(V)$ 是所有函数 $V\longrightarrow\mathbb C$ 组成的 $\mathbb C$-代数. 这样的函数称为多项式函数，如果它是 $\mathbb C^n$ 上某个多项式函数的限制. 因而，函数的限制映射给出满同态
\begin{align*}
\mathcal O(\mathbb C^n)=\mathbb C[X_1,\ldots,X_n]\twoheadrightarrow\mathcal O(V).
\end{align*}
设 $I$ 是其核，即所有在 $V$ 上为零的函数 $f\in\mathcal O(\mathbb C^n)$ 组成的理想. 根据最后一个推论，$I$ 有限生成，由一些多项式 $f_1,\ldots,f_r\in\mathbb C[X_1,\ldots,X_n]$ 生成. 因而
\footnote{译注：原文此式末尾没有写出 $=0$；按上下文，此处应要求所有 $f_i$ 的值均为零.}
\begin{align*}
V=V(I)&:=V(f_1,\ldots,f_r)\\
&:=\{(z_1,\ldots,z_n)\in\mathbb C^n,\ f_1(z_1,\ldots,z_n)=\cdots=f_r(z_1,\ldots,z_n)\}.
\end{align*}
[事实上，包含关系 $V\subset V(I)$ 是显然的；又因为 $V$ 是代数集，因而可写成 $V(f'_1,\ldots,f'_{r'})$，其中 $f'_i$ 是另一些多项式，所以对每个 $i$ 都有 $f'_i\in I$，从而 $V(I)\subset V(f'_1,\ldots,f'_{r'})=V$.] 因此可见，商 $\mathbb C[X_1,\ldots,X_n]/I=\mathcal O(V)$ 的泛性质给出双射
\footnote{译注：原文此处印作 $z\in(f\mapsto f(z))$；作为映射的对应规则，应读作 $z\mapsto(f\mapsto f(z))$.}
\begin{align*}
V\xrightarrow{\sim}\operatorname{Hom}_{\mathbb C\text{-}\mathrm{alg}}(\mathcal O(V),\mathbb C),\qquad z\in(f\longmapsto f(z)).
\end{align*}
再一次，空间 $V$ 的点可以解释为其多项式函数 $\mathbb C$-代数上的求值同态.

此外，我们已经看到，Noether 环的任何商环都是 Noether 环. 结合前一个推论可知，$\mathcal O(V)$ 是 Noether 环（但未必是整环！）.

\subsection{Laurent 多项式}\label{sec:galois-1-4-6}
这是 $\mathcal N=\mathbb Z$ 的例子. 因而 $A$-代数 $A[\mathbb Z]$ 有一组基 $(e_n)_{n\in\mathbb Z}$，对所有 $n,m\in\mathbb Z$ 都满足 $e_ne_m=e_{n+m}$. 令 $X:=e_1$. 对每个 $n\in\mathbb Z$ 都有 $e_n=X^n$，而 $A[\mathbb Z]$ 的每个元素都唯一地写成 $f=\sum_{n\in\mathbb Z}a_nX^n$.
\begin{definition*}
代数 $A[\mathbb Z]$ 称为``\newterm{galois-laurent-polynomial}{Laurent 多项式}{Laurent polynomial}代数''. 通常记作 $A[X,X^{-1}]$，有时也记作 $A[X^{\pm1}]$.
\end{definition*}
$A[X,X^{-1}]$ 的泛性质表述如下：
\begin{proposition*}
对任意配备了一个可逆元素 $b$ 的 $A$-代数 $B$，存在唯一的 $A$-代数同态 $\varphi\colon A[X,X^{-1}]\longrightarrow B$，使得 $\varphi(X)=b$. 换言之，对任意 $A$-代数 $B$，映射 $\varphi\longmapsto\varphi(X)$ 诱导双射
\begin{align*}
\operatorname{Hom}_{A\text{-}\mathrm{alg}}(A[X,X^{-1}],B)\xrightarrow{\sim}B^\times.
\end{align*}
\end{proposition*}
可以改述为：给定同态 $A[X,X^{-1}]\longrightarrow B$，就相当于给定 $B$ 的一个可逆元素，即 $X$ 的像.
\begin{proposition*}
若 $A$ 是整环（相应地，Noether 环），则 $A[X,X^{-1}]$ 也是整环（相应地，Noether 环）.
\end{proposition*}
\begin{proof}
设 $A$ 是整环. ``通常的''多项式环 $A[X]$ 包含在 $A[X,X^{-1}]$ 中，而且我们知道它是整环. 设 $f,g\in A[X,X^{-1}]$ 满足 $fg=0$. 存在整数 $n\in\mathbb N$，使得 $fX^n\in A[X]$ 且 $gX^n\in A[X]$. 于是，在 $A[X]$ 中成立的等式 $(fX^n)(gX^n)=0$ 蕴含 $fX^n=0$ 或 $gX^n=0$. 由于 $X$ 在 $A[X,X^{-1}]$ 中可逆，故 $f=0$ 或 $g=0$. 因而 $A[X,X^{-1}]$ 是整环.

再设 $A$ 是 Noether 环. 我们要把 $A[X,X^{-1}]$ 表示为一个已知的 Noether 环的商环. 为此，考虑将 $X$ 映为 $X$、将 $Y$ 映为 $X^{-1}$ 的唯一同态 $A[X,Y]\longrightarrow A[X,X^{-1}]$（由泛性质给出）. 它也将 $X^n$ 映为 $X^n$，将 $Y^n$ 映为 $X^{-n}$. 因而可见，它是满射，因为其像包含 $A[X,X^{-1}]$ 的一组基. 我们已经看到 $A[X,Y]$ 是 Noether 环，故 $A[X,X^{-1}]$ 也是 Noether 环.
\end{proof}
\begin{exercise*}
证明：将 $X$ 映为 $X$、将 $Y$ 映为 $X^{-1}$ 的同态 $A[X,Y]\longrightarrow A[X,X^{-1}]$ 的核，是由 $XY-1$ 生成的理想，从而
\begin{align*}
A[X,X^{-1}]=A[X,Y]/(XY-1).
\end{align*}
\end{exercise*}

\subsection{$n$ 个未定元的 Laurent 多项式}\label{sec:galois-1-4-7}
这是 $\mathcal N=\mathbb Z^n$ 的例子. 代数 $A[\mathbb Z^n]$ 包含子代数 $A[\mathbb N^n]$，因而包含前面定义的元素 $X_1,\ldots,X_n$. 可见，Laurent 单项式 $X_1^{\nu_1}\cdots X_n^{\nu_n}$（其中 $\nu=(\nu_1,\ldots,\nu_n)\in\mathbb Z^n$）构成 $A[\mathbb Z^n]$ 的一组 $A$-基. 这个环也记作 $A[X_1^{\pm1},\ldots,X_n^{\pm1}]$. 给定同态 $A[X_1^{\pm1},\ldots,X_n^{\pm1}]\longrightarrow B$，等价于给定由 $B$ 的可逆元素组成的 $n$ 元组 $(b_1,\ldots,b_n)$. 如同普通多项式的情形一样，由此可见
\begin{align*}
A[X_1^{\pm1},\ldots,X_n^{\pm1}]\simeq A[X_1^{\pm1},\ldots,X_{n-1}^{\pm1}][X_n^{\pm1}].
\end{align*}
由归纳法可知，若 $A$ 是整环（相应地，Noether 环），则 $A[X_1^{\pm1},\ldots,X_n^{\pm1}]$ 也是如此.
\begin{remark*}
一般而言，$A$ 为整环并不蕴含 $A[\mathcal N]$ 为整环. 例如，当 $\mathcal N=\mathbb Z/2\mathbb Z$ 时，有 $A[\mathbb Z/2\mathbb Z]\simeq A[X]/(X^2-1)$（练习）. 在其中，$(X-1)(X+1)=0$，但 $X-1$ 和 $X+1$ 都非零. 至于``$A$ 为 Noether 环蕴含 $A[\mathcal N]$ 为 Noether 环''这一性质，无限多个变量的多项式 $A[\mathbb N^{\mathbb N}]=A[X_1,X_2,\ldots]$ 的例子说明，它并不总是成立.\footnote{译注：原文确实印作 $\mathbb N^{\mathbb N}$. 按前文的幺半群代数构造，无限变量多项式对应的指数幺半群应为有限支撑序列组成的 $\mathbb N^{(\mathbb N)}$，而不是全部序列组成的 $\mathbb N^{\mathbb N}$. 此处保留原式并指出这一疑点.} 不过，如果幺半群 $\mathcal N$ 由有限多个元素生成，这一性质就成立（练习）.
\end{remark*}



